\section{Counterfactual Forecasting Framework}\label{sec:framework}

\subsection{Predictive Counterfactual Formulation}

The forecasting task is defined over district-level time series. For Ridge and LightGBM, a separate model is estimated for each pollutant $p$ using autoregressive windows that combine the recent history of that pollutant with district traffic, meteorology, and cyclic calendar features. If $x_t$ denotes the exogenous and calendar features at time $t$ and $y_t^{(p)}$ the concentration of pollutant $p$, the pollutant-specific forecast has the generic form
\[
\hat{y}_t^{(p)} = f_{\theta}^{(p)}(y_{t-w}^{(p)}, \ldots, y_{t-1}^{(p)}, x_{t-w}, \ldots, x_{t-1}).
\]

The stacked LSTM is instead estimated as a multivariate sequence model. It uses windows of exogenous and calendar inputs and predicts the pollutant vector jointly:
\[
\hat{\mathbf{y}}_t = g_{\theta}(x_{t-w}, \ldots, x_{t-1}),
\]
where $\hat{\mathbf{y}}_t$ contains the five pollutant forecasts at time $t$. At post-policy times, the counterfactual predictions represent the pollution levels expected under continuation of the pre-policy regime, whereas the realized observations reflect the post-policy setting.

The interpretation is intentionally modest. Positive gaps of the form $\hat{y}_t - y_t > 0$ are treated as evidence consistent with lower observed post-policy concentrations relative to the no-change forecast. They are not treated as proof of strong causal identification, because the model-based counterfactual remains sensitive to predictive misspecification and to the information encoded in the pre-policy period.

The estimand also matters for how predictions are generated after the boundary. For the autoregressive Ridge and LightGBM models, a recursive rollout is required if the forecast is meant to represent continuation of the pre-policy regime. Feeding observed post-policy pollutant lags back into the model would update the autoregressive state with treated outcomes and would therefore answer a different question. A direct one-step evaluation based on observed post-policy lags is still informative as a symmetry check for predictive ranking, but not as the primary no-change counterfactual.

\subsection{Preprocessing}

The preprocessing pipeline consists of district assignment for sensors, temporal alignment of air-quality, meteorological, and traffic records, hourly reindexing over 2021--2023, interpolation of missing exogenous values, and feature scaling where required by the model family. An IQR-based outlier treatment is retained as a clipping rule applied to pre-policy ranges so that the hourly sequence remains continuous for forecasting. Rather than treating that rule as unquestioned preprocessing, the empirical design evaluates its removal as a sensitivity case.

To keep the paper concise, generic textbook descriptions of forecasting metrics are omitted. What matters here is the sequence of design choices that can plausibly change the sign or magnitude of the post-policy comparison.

\subsection{Temporal Windows}

Temporal dependence is tested with two explicit windows in the main design: $w=6$ and $w=24$. These choices represent short-term and day-scale context, respectively, and are long enough to test whether additional temporal memory changes the inferred post-policy signal while keeping the comparison compact and interpretable.

Using only two windows in the core tables is a deliberate design choice. The objective is not an exhaustive hyperparameter search, but a targeted robustness check. If the sign of the policy-consistent effect changes materially between $w=6$ and $w=24$, then the substantive interpretation becomes more fragile. If it remains directionally stable, the result is easier to defend. A supplementary rolling-origin sensitivity later introduces $w=12$ only to verify that the preferred short-range window is not an artifact of the main split.

\subsection{Compared Models}

The extended manuscript compares three model families chosen to cover distinct inductive biases without expanding the paper beyond its editorial scope.

\paragraph{\textbf{Ridge}}
Ridge regression is the linear benchmark. It provides a transparent baseline and helps determine whether the post-policy signal already appears under a restricted linear specification.

\paragraph{\textbf{LightGBM}}
LightGBM represents a strong non-linear machine-learning baseline for tabular data. In this context it captures feature interactions and non-linearities without requiring a heavier architectural commitment. The model is fitted on flattened historical windows, which makes it directly comparable to the linear baseline while preserving the same temporal context.

\paragraph{\textbf{LSTM}}
The LSTM model provides the sequential deep-learning alternative \cite{10.1162/neco.1997.9.8.1735}. In the implemented comparison, it is specified as a stacked five-layer multivariate network with dropout after each LSTM block and joint prediction of the five pollutants. This design preserves a recognizably sequential architecture while keeping the model comparison manageable.

The model set is intentionally limited to these three families. Adding graph neural networks or transformers would increase implementation cost and discussion burden without clearly improving the interpretability of a compact policy-oriented comparison. The available data benefit more from controlled robustness than from architectural proliferation.

\subsection{Effect Summaries}

Post-policy behavior is summarized first through the positive impact percentage (PIP), defined as the share of post-policy observations for which the counterfactual prediction exceeds the realized value:
\[
\mathrm{PIP} = \left( \frac{1}{N} \sum_{i=1}^{N} \mathbf{1}(\hat{y}_i > y_i) \right) \times 100.
\]

PIP is supplemented with two signed residual summaries:
\[
\mathrm{MR} = \frac{1}{N} \sum_{i=1}^{N} (\hat{y}_i - y_i),
\qquad
\mathrm{MdR} = \mathrm{median}_{i=1,\ldots,N} (\hat{y}_i - y_i).
\]
Here $\mathrm{MR}$ is the mean residual and $\mathrm{MdR}$ is the median residual over the post-policy period. Positive values indicate that observed pollution is lower than the no-change forecast, whereas negative values indicate the opposite. Using all three summaries gives a more informative picture than PIP alone because it separates directional consistency from effect magnitude and reduces dependence on outliers.
